
@@M0@@ (VI 5.5), on
peut supposer que @@M1@@ est un faisceau inversible. On a alors : @@M2@@ ; donc @@M3@@.

Les @@M4@@, @@M5@@ sont nuls pour @@M6@@. Il suffit donc de montrer que @@M7@@. Or @@M8@@ est la classe dans @@M9@@ de @@M10@@, et @@M11@@ celle de @@M12@@. L'égalité :
@@M13@@
donne :
@@M14@@
on en déduit le résultat.

\textbf{Théorème 3.6. (Riemann--Roch).} \uline{Soit} @@M15@@ \uline{un morphisme projectif d'intersection complète} (1.1), \uline{et dont la dimension relative virtuelle} @@M16@@ \uline{est constante} (1.9). \uline{On suppose que} @@M17@@ \uline{est quasi-compact}, \uline{et admet un faisceau ample}. \uline{Avec les notations de} 3.3, \uline{on a alors pour tout} @@M18@@ \uline{la relation}:
@@M19@@
@@M20@@ \uline{désignant l'élément dual de la classe du complexe cotangent relatif} @@M21@@ (2.5 \uline{et} 2.3).

En d'autres termes, le diagramme
@@M22@@
est commutatif.

